\section{Discussion}


This study examined whether and how media coverage of the Republican and Democrats evolves in relation to the issue of Ukraine Aid from March 2022 to December 2024, treating the rapid transformation of once-bipartisan foreign policy issue as a naturally occurring case for observing party brand. Drawing insights from \citet{riegerTheyShouldntReceive2025}, we tracked polarization and within-party coherence through the General Ukraine Aid Stance dimension and three explanatory dimension: 1) how Ukraine is characterized, 2) which external threat should be prioritized, and 3) whether foreign commitment overweights domestic priorities.

We found that General Ukraine Aid Stance polarization increased substantially over the 34 months, with a Kendall's $\tau$ = .330 (p = .006), and a Sen slope of 0.29 points per month. We then compared three descriptions of the temporal trajectory of polarization: a single gradual trend, a set of persistent shifts, and a model with sets of shifts and a common trend running through them. The shift-only model achieved the lowest BIC, with the common trend contributing no significantly additional information to the estimation once shifts were admitted. The selected model placed three regime onset month to September 2022, April 2023, and September 2023, leaving four regimes with polarization levels of 28.2, 37.27, 28.74, and 39.91. Polarization increased substantively first, but returned to its start point, then increased again and sustained for 16 months. The evidence thus supports a temporal structure that contains multiple punctuated shifts rather than a gradual trend.

Both parties moved significantly towards opposition on the issue of Ukraine Aid, though Republicans moved further. On the Sen slope, Republican moved on average 0.42 points towards the opposition per month, and Democrats moved roughly 0.17 points per month. The median difference between their growth is 0.29 points per month, also significant. This finding is consistent with the notion that minority party are incentivized and rewarded more to establish differentiation on emerging issues. Among the three explanatory dimension, Resource Allocation showed the largest contemporaneous contribution to the General Ukraind Aid Stance polarization, accounting for 61.2\% of modeled R2, superseding Ukraine Characterization (24.3\%) and Threat Priority (14.4\%). However, this ordering did not survive the stationary bootstrap, therefore, we cannot confirm reliably, within our corpus sample, that this finding generalizes.

The Party Brand Index increased significantly across the 34 months, with a Sen slope of 0.36 points per month. The BIC selected model had one onset month in September 2023, leaving two regimes averaging at 53.72 and 61.30 index points. However, the underlying General Aid Stance polarization and within-party coherence did not move together. Neither party's polarization coefficient was reliably different from zero when predicting its own within-party coherence of the same month. What did change was Republican coherence: coverage about Republican converged by a fitted 25.3 IQR points over the studied period, while Democrat's coherence showed a non-significant decline. Therefore, at least for Republican, it has both differentiated sharply from Democrat on the issue of Ukraine aid, and has, over the course of 34 months, consolidated their internal messages regarding their stance.

Overall, our findings show that a foreign policy issue can acquire a partisan brand within months, and that the process was neither smooth, nor symmetric, nor synchronized. It was not smooth. Both polarization and the Party Brand Index moved between persistent levels rather than gradual trend, and both settled into their final regime in September 2023.

It was not symmetric. Both parties moved toward opposition, but Republicans moved further, so the widening polarization reflected a difference in the extent of directional movement rather than one party moving against a fixed position of another.

It was not synchronized. Both conditions a party brand requires advanced for the Republican Party, which moved further from the Democratic position and consolidated on a more consistent position of their own, yet the two did not move together from month to month.

\subsection{Theoretical Implications for Party Brand}

Party brand theory holds that parties are motivated both to occupy distinct positions and to maintain consistent messaging, so that voters can use party labels as informational shortcuts \citep{lupuPartyBrandsPartisanship2013,lupuPartyBrandsCrisis2016}. Across the 34 months, both conditions advanced for the Republican Party. Its represented position moved further from the Democratic position, and its coverage converged on a more consistent position of its own. The content of the theory is therefore upheld in the case we examined. What our findings adds further is its timing. Polarization and within-party coherence did not move together from month to month.

There instead seems to be a lag for within-party coherence, and there are few explanations as to why. First, polarization and consolidation are not equally costly. A party's represented position can move as soon as a few prominent voices push out an extreme view, whereas its represented range narrows only if most of its voices follow. The first requires position-taking by individuals; the second requires the remainder of the party to fall into line, which is the work of coordination and resolution of intraparty disputes. Second, our two measures track different parts of the same distribution. Polarization is a difference in party means and responds to the tails, while the IQR is constructed not to. A shift driven by a vocal minority therefore registers in one measure before the other, so temporal misalignment between them might be expected rather than anomalous.

A third possibility concerns how news coverage represents parties rather than how parties behave. Our Literature Review established that perceived brands form through indirect communication, and that argument has a consequence we can now state. When a party is repositioning, its internal disagreement is itself newsworthy, so coverage may amplify the range of intraparty opinion at precisely the moments when the party is trying to narrow it. If that holds, represented coherence is structurally harder to achieve than represented distinctiveness, and the lag we observe would be a property of the medium rather than of the party. Our design cannot separate this reading from the previous two.

Democratic coverage shows the complement to the Republican pattern. Democrats held an upward but more maintained position, with a coherence that did not narrow and in fact widened slightly without reaching significance. Republicans, holding the minority position in the Senate and facing unified Democratic control until late 2022, had reason to establish ownership on an emerging issue \citep{abou-chadiModerateNecessaryRole2016,kollbergChallengerAdvantageHow2025}. Democrats faced the costs of abandoning a signature policy of the incumbent administration. What our findings add is that the challenger's project extended beyond positional movement to internal consolidation, while the incumbent's did neither.

Additionally, this asymmetry is visible only because we measured both conditions at once. Party brand theory names distinctiveness and consistency as jointly necessary, but empirical work has largely operationalized the first, and a measure of divergence cannot register whether the second is keeping pace. Constructing the Party Brand Index as a geometric mean makes that question answerable, because a conjunctive measure is dampened whenever its components develop unevenly. The index increased across the period and stepped up once, in September 2023, on a base where one party consolidated and the other did not. Its rise is therefore conservative with respect to what joint development would have produced, and the gap between the two is itself a point of interest. Therefore, the index does more than summarize brand conditions. It provides a way of asking whether a brand is developing on both conditions or on one, which is a question the theory poses and existing measures cannot answer.

\subsection{Foreign Policy Polarization}

The Ukraine aid case also speaks to longstanding debates about whether politics stops at the water's edge. \citet{bryanPrevalenceBipartisanshipUS2022} find that extreme partisan division on foreign policy remains the exception and that international issues are still less polarized than domestic ones. Our findings do not contradict that general claim, but they show how quickly a particular issue can move. An issue that began with broad elite agreement and little established partisan disagreement acquired a persistent partisan structure within eighteen months. \citet{guisingerMappingBoundariesElite2017} show that the influence of party cues varies with how much baseline disagreement an issue already carries, which implies that an issue arriving without such a baseline is not thereby protected from acquiring one. Partisan sorting, media fragmentation, and intensified electoral competition may now expose foreign policy issues to the same dynamics that characterize domestic politics.

The dimensional result speaks to how that exposure operates. \citet{jeongDivisionWatersEdge2019} argue that foreign policy polarization is largely derivative of broader domestic polarization, and \citet{borgBattleSoulThis2024} extends this to competing visions of American identity and priorities. Of the three explanatory dimensions, the framing that set foreign commitments against domestic priorities showed the closest contemporaneous alignment with the aid stance gap, contributing 0.18 of the complete model R2 of 0.294. This is the ordering the spillover account anticipates: connecting a novel foreign policy issue to well-established domestic issue appears to be the route by which it polarized, rather than disagreement about Ukraine itself or about which external threat deserves priority. Both simultaneous contrasts included zero and Resource Allocation ranked first in only 50.5 percent of dependence-preserving replicates, so the ordering is descriptive rather than established.

Our findings also speak about how apparent consensus should be interpreted. \citet{friedrichsPolarizationUSForeign2022} note that interparty agreement can coexist with intraparty division, and that aggregate party differences can obscure what is happening inside each party. Our two measures separate these. The between-party polarization widened while Republican's coherence narrowed and Democratic coherence did not, which is a pattern no measure of polarization alone would reveal. This suggests that studies reporting foreign policy polarization from party means may be describing only half of what changed.

Finally, our findings do not imply that all foreign policy issues will polarize. Ukraine aid may represent a particular type of issue that is novel, initially unsettled, and open to reframing around domestic priorities.

\subsection{The Structure of the Polarization Trajectory}

The joint model located three estimated regime onsets, in September 2022, April 2023, and September 2023. The trajectory rose, returned to approximately its starting level, then rose again to a level it held for the final 16 months. This structure matters beyond model fit. Party brand theory describes differentiation as a strategic project, and a project that proceeds in stages looks different from one that accumulates. Our trajectory shows a period of elevated branding that did not hold, followed by one that did. Whichever mechanism produced the durable shift, it was not simply more of what produced the first.

For the polarization to grow, one party has to move away from the other, and around the first onset both moved. On the Republican side, opposition shifted from the party's fringe to its leadership rank. Through the first half of 2022 the position belonged to a small group of members, but in October 2022 then Minority Leader Kevin McCarthy told The Hill that a Republican majority would not write a blank check for Ukraine \citep{brooksMcCarthyWarnsNo2022}. On the Democratic side, the one move in the opposite direction was suppressed within a day. Thirty members of the Congressional Progressive Caucus released a letter urging the White House to pursue direct diplomacy with Russia on 24 October 2022, and the caucus withdrew it the following afternoon, with its chair stating that Democrats were united in their commitment to Ukraine and that the letter was being conflated with McCarthy's position \citep{pilkingtonProgressiveDemocratsRetract2022}. Democratic elites objected that the letter undercut a position of support at precisely the moment Republicans were divided over further aid \citep{rajuJayapalWithdrawsLetter2022}. Both movements are what the asymmetric incentives reviewed would predict: the challenger attempts to own the emerging issue as an election approach, and the incumbent seeks to maintain a position it already owns. The regime then held for seven months, perhaps due to the conversion of a campaign position into a governing one. Republicans took the House majority in January 2023, so a commitment McCarthy made while seeking the speakership became the position of the chamber that controls appropriations. What had been a claim about what a majority would do became a constraint on what the majority did.

The second onset returns polarization to approximately its opening level. The most prominent Republican statement against aid in this window is from within the party. Ron DeSantis described the war as a territorial dispute outside vital United States interests in March 2023 \citep{garrityRepublicansSlamDeSantis2023}, and senior Senate Republicans including Mitch McConnell, Marco Rubio, Lindsey Graham and Thom Tillis publicly rejected that characterization \citep{garrityRepublicansSlamDeSantis2023}. A party mean is pulled by whichever voices are covered, so a month in which prominent Republicans defend aid loudly will narrow the gap even as the party grows more internally divided. The fall during this regime is therefore compatible with Republicans becoming less unified rather than more supportive. There are also legislative agenda pushing for de-fund. The debt ceiling crisis dominated Congress through spring 2023, and legislators described Ukraine funding as delayed and on the back burner while it ran \citep{dressUkraineAidDries2023}. The anticipated counteroffensive, which began in June, gave a further reason to defer the next funding conversation. Ukraine aid was not settled during these five months, it was simply not the main portion of the contestation, and a position that is not being contested is not being differentiated. This regime is also the shortest and the least securely located. The six-month minimum duration moves its onset from April to March, which is the month of the DeSantis remark and Republican conflict, so our data cannot separate the two interpretations.

The third onset produces the highest regime and the only one that holds. September 2023 is when opposition to Ukraine aid acquired legislative debate. Zelensky addressed senators at the Capitol on 21 September, and McCarthy declined to grant him a joint address \citep{grovesGovernmentFundingPackageOct2023}. In the same period, close to half of House Republicans voted to strip \$300 million for Ukraine from a defense spending bill \citep{frekingUkraineAidDropped2023}. On 30 September, the government shutdown was only resolved after \$6 billion for Ukraine was removed \citep{harrisCongressAvoidsGovernment2023}, and McCarthy was removed from the speakership three days later \citep{mascaroHow6Billion2023}. In 2022 the Republican position was a campaign commitment that a majority had to honor. In 2023, some elites made the position to de-fund Ukraine binding on the party. The level established here held for the final sixteen months, through the primary season and the general election campaign.

The Party Brand Index has a single onset, in September 2023, against three in polarization, and the difference perhaps lies in the internal chaos of Republican. As illustrated above, Republican did not necessarily gain much coherence in its more prominent position to de-fund Ukraine after September 2023. In September 2023, polarization reached its highest sustained level, Republican coherence reached its best, and Democratic coherence recovered from its low. It marks the one moment in the study period when both conditions a party brand requires were moving in the same direction at once.

\subsection{Limitations}

First, the source of measurement limits our finding to how parties were represented in news coverage rather than what elites said directly. Journalists select which statements to report, which voices to amplify, and how to frame partisan conflict, and if that selection favors conflictual statements, our measures may overstate represented polarization relative to the underlying positions. We argued that for that exact reason, news is the appropriate source of measurement, because perceived party brands form through indirect communication as well as direct communication. \citet{nomikosAmericanSocialMedia2025} also found that the same trend through Twitter data, where liberal Twitter users expressed greater support for Ukraine than conservatives did. It suggests that at least on the aggregate level, what we recovered is not an artifact of journalistic selection alone. The choice to use news coverage also means that we cannot separate party behaviors from the media representation. The lag we infer between differentiation and consolidation is consistent with parties taking time to bring their members into line, and it is equally consistent with coverage amplifying internal disagreement precisely when a party repositions. Distinguishing these requires comparing the same statements as delivered and as reported, which our corpus does not contain.

Second, we examined one foreign policy issue, in one national context, during one historical period. This design enabled detailed temporal tracking but cannot establish whether the observed dynamics would generalize to other foreign policy debates. Ukraine aid may be a particular type of issue: novel, initially bipartisan, then opened to contestation with domestic priorities. The rapid polarization we documented could reflect something specific to Ukraine, or it could reflect broader changes in the U.S. political environment that would expose any foreign policy issue to similar change.

Third, several features of our annotation suite limit what the measures can support. We ran a single prompt rather than a set of variants, so we did not conduct a prompt stability analysis that \citet{tornbergBestPracticesText2024} argued a standard for LLM-based annotation. We cannot therefore report how far our estimates would move under alternative wordings of the same instructions. We used a single model, and exact reproduction requires access to gpt-5.1-2025-11-13. Our human reference sample is small, especially for explanatory dimensions. Agreement on NA assignment is the weakest part: human coders disagreed about applicability on 1 percent of ratings while the model disagreed with human on roughly one in six, and because NA determines which sentences enter each dimension, this affects coverage more than it affects the scores. Resampling the 50-sentence reference set 5,000 times and recomputing the three criteria, all three clear in 74\% of replicates, so a different draw of sentences would have produced a different verdict in roughly one case in four. Comparing the validation annotation with the production annotation of the same 680 sentences, 225 differed on at least one dimension, with exact agreement ranging from 78.7 percent on General Ukraine Aid Stance to 94.7 percent on Threat Priority. The model is therefore highly stable when the same sentences are scored together and less so when they are scored again on a later occasion. Taken together, those validation statistics support good but limited sample-specific agreement between the model and human coders on this corpus.

\subsection{Wider Implications and Future Directions}

These findings carry implications for understanding how foreign policy issues become domesticated into partisan conflict. The speed with which Ukraine aid acquired a partisan structure should warn practitioners against assuming sustained consensus. Our analysis shows polarization moving between persistent levels rather than drifting, so a period of apparent stability offers no assurance that the prevailing level will hold. For policymakers seeking to maintain cross-party coalitions, framing appears to matter: the dimension tying Ukraine aid to domestic spending showed the closest alignment with the aid stance gap, which suggests that strategies insulating foreign policy questions from domestic budgetary framing may sustain agreement longer.

The design we employed offers a template for tracking party brand formation on emerging issues. Annotating individual sentences along several predetermined dimensions, aggregating to articles so that no single article contributes disproportionately, and then aggregating to months preserves dimensional specificity while allowing measurement at a scale that repeated survey collection could not reach. This approach could be applied to other issues undergoing partisan contestation, such as climate policy, immigration, or technology regulation, to identify which framings move with polarization and when the structure of a debate changes.

A direct test of our timing account is available to future work. If consolidation follows differentiation rather than accompanying it, a party's internal range should narrow at some lag behind movement in the between-party gap, and that lag should be estimable. We could not estimate it here. With 34 monthly observations, a distributed-lag specification spends degrees of freedom faster than it recovers signal, and the lag length would have to be selected from the same data used to fit it. A series spanning several issue cycles would support both steps. Such a design would also allow the two components to be separated by what drives them, since between-party divergence can be produced by a small number of prominent voices while within-party coherence requires the center of a party's coverage to move.

Three further extensions would strengthen these findings. Comparative analysis across multiple foreign policy issues would establish whether Ukraine aid represents a unique case or exemplifies broader dynamics. Comparison of mediated discourse with direct elite speech, such as floor statements or campaign communications, would quantify how journalistic filtering shapes polarization estimates and would address the third limitation above. And linking represented elite polarization to mass opinion dynamics would test whether the rising Party Brand Index corresponds to clearer public understanding of party positions and to stronger partisan attachment, which is the validation the index most needs.
